% Supplementary S15 — Identifiability checks
% 2026-06-05
% Candidates reported: S08-robust (6,−42), S09-primary (2,+24).
% S08-stable excluded (dropped at selection stage).
% All tests use PCA-basis voxel synthesis (spatial PCA(20) on SRM-projected amplitudes).

\subsection*{\textbf{S15. Identifiability checks}}
\label{app:identifiability}

Four pre-specified checks assessed whether the production argmin
is identifiable (Test~1), what the algorithm's noise floor is
(Test~2a), and whether the fit is specific to CVD data
(Tests~2b--2c).
Each test was run on the two production candidates (S08-robust,
S09-primary).
All tests use PCA-basis voxel synthesis (spatial PCA(20) on
SRM-projected amplitudes) as the feature space.
FDR correction (Benjamini--Hochberg, $\alpha = 0.05$) was applied
across 6 tests (2 candidates $\times$ 3 main tests); 0/6 reached
significance.

% ─────────────────────────────────────────────────────────────────────────────
\paragraph{Test 1 — Voxel-level parameter recovery.}
% ─────────────────────────────────────────────────────────────────────────────

Synthetic fMRI data were generated at the production argmin
(PCA-basis; 7~HC donors $\times$ 20 noise draws $= 140$ samples
per candidate) with GT-consistent JND: psychophysical thresholds were
scaled to match the distortion magnitude at the ground-truth
$(\beta_s, \beta_c)$.
The same fitting pipeline was re-run on each synthetic sample.
Pass criterion: fraction within $10^\circ$ on both axes
($f_{10^\circ}$) $\geq 0.5$ and $|$bias$| < 10^\circ$ on both axes.

\begin{table}[h]
\centering
\begin{tabular}{lccc}
\hline
Candidate & bias $(\beta_s,\; \beta_c)$ & $f_{10^\circ}$ & Verdict \\
\hline
S08-robust  $(6^\circ, -42^\circ)$  & $(+16^\circ,\; -4.7^\circ)$  & 0.26 & FAIL \\
S09-primary $(2^\circ, +24^\circ)$  & $(+11^\circ,\; -27^\circ)$  & 0.14 & FAIL \\
\hline
\end{tabular}
\caption{Test~1 results. Both production candidates fail the recovery
criterion ($f_{10^\circ} < 0.5$).
The non-dominant axes ($|\hat\beta_s| \leq 6^\circ$ for both
participants) fall below the effective parameter uncertainty
established by Test~2a and are not recoverable.
For S08-robust, the dominant axis ($|\hat\beta_c| = 42^\circ$, bias
$= 4.7^\circ$) shows lower bias than the non-dominant axis
($16^\circ$), consistent with partial magnitude recoverability
where the axis magnitude exceeds the noise floor.}
\end{table}

% ─────────────────────────────────────────────────────────────────────────────
\paragraph{Test 2a — Origin-recovery null.}
% ─────────────────────────────────────────────────────────────────────────────

Synthetic fMRI data were generated at GT $= (0, 0)$ using real HC
JND (no synthetic psychophysical signal; PCA-basis).
At the null, real HC JND carries no CVD-consistent signal, making
this a noise floor measurement free of synthesis design
contamination.
Pass criterion: $|$bias$| < 5^\circ$ on both axes.

\begin{table}[h]
\centering
\begin{tabular}{lcccc}
\hline
Candidate
  & $|\hat\beta_s|_{\rm med}$ (IQR)
  & $|\hat\beta_c|_{\rm med}$ (IQR)
  & $\beta_c$ p95
  & $f_{10^\circ}^{\rm origin}$ \\
\hline
S08-robust  & $22^\circ\;(40)$ & $26^\circ\;(10.5)$ & $44^\circ$ & 0.00 \\
S09-primary & $16^\circ\;(17.5)$ & $24^\circ\;(9)$  & $48^\circ$ & 0.00 \\
\hline
\end{tabular}
\caption{Test~2a results. $f_{10^\circ}^{\rm origin} = 0/140$ for
both candidates: no sample was recovered within $10^\circ$ of the
origin under a null ground truth.
Effective parameter uncertainty (PCA-basis): $\sim\!20^\circ$ ($\beta_s$)
and $\sim\!25^\circ$ ($\beta_c$).}
\end{table}

% ─────────────────────────────────────────────────────────────────────────────
\paragraph{Test 2b — HC pseudo-CVD specificity.}
% ─────────────────────────────────────────────────────────────────────────────

Each HC participant's real amplitude data were substituted as a
fake CVD participant (7 carriers per candidate, deterministic).
The production CVD fit distance from the origin was ranked among
the 7 HC pseudo-CVD distances (rank$_{\rm dist}$, one-sided high).
Pass criterion: rank $= 1.0$ (CVD fit exceeds all HC distances).
This test is descriptive only and was not used as a selection
criterion.

\begin{table}[h]
\centering
\begin{tabular}{lcc}
\hline
Candidate & rank$_{\rm dist}$ & Verdict \\
\hline
S08-robust  & 0.875 & FAIL \\
S09-primary & 0.875 & FAIL \\
\hline
\end{tabular}
\caption{Test~2b results. Neither candidate exceeds all HC
pseudo-CVD distances.}
\end{table}

% ─────────────────────────────────────────────────────────────────────────────
\paragraph{Test 2c — Color-label permutation.}
% ─────────────────────────────────────────────────────────────────────────────

CVD trial labels were randomly shuffled ($N = 1000$ permutations);
HC data were held fixed.
Pass criterion: $p_{\rm perm} < 0.05$ (one-sided, lower = better).

\begin{table}[h]
\centering
\begin{tabular}{lccc}
\hline
Candidate & Real loss & Perm 5\% cut & $p_{\rm perm}$ \\
\hline
S08-robust  & $-2.892$ & $-3.136$ & 0.167 \\
S09-primary & $-1.681$ & $-3.053$ & 0.471 \\
\hline
\end{tabular}
\caption{Test~2c results. Neither candidate reaches significance
against the label-permutation null.}
\end{table}

% ─────────────────────────────────────────────────────────────────────────────
\paragraph{FDR summary and basis caveat.}
% ─────────────────────────────────────────────────────────────────────────────

After FDR correction across 6 tests (2 candidates $\times$ 3 main
tests; BH $\alpha = 0.05$), no test reaches significance.
The production argmin is treated as a descriptive embedding of the
CVD-specific distortion direction throughout.

All results above use the PCA-basis loss.
For the protan participant, the sign of $\hat\beta_c$ (mechanism class) was confirmed
under the PCA-basis loss only.
The SRM-basis loss may yield a different mechanism class for this
participant; sign stability under the SRM basis was not verified.
